\documentclass[aspectratio=169,10pt]{beamer}
% ═══════════════════════════════════════════════════════════════════════════════
% SHARED PREAMBLE — Foundation Models Course (HIT)
% To be \input by each lecture file AFTER \documentclass
% ═══════════════════════════════════════════════════════════════════════════════

% ─────────────────────────────────────────────────────────────────────────────
% BEAMER THEME & BASIC SETUP
% ─────────────────────────────────────────────────────────────────────────────
\usetheme{Madrid}
\usecolortheme{default}

% Remove navigation symbols
\beamertemplatenavigationsymbolsempty

% ─────────────────────────────────────────────────────────────────────────────
% COLORS — HIT Branding
% ─────────────────────────────────────────────────────────────────────────────
\definecolor{hitblue}{RGB}{0,51,102}        % Primary: HIT Navy
\definecolor{hitgold}{RGB}{192,148,0}       % Accent: HIT Gold
\definecolor{hitlight}{RGB}{230,240,250}    % Light background

% Override Beamer palette
\setbeamercolor{primary}{fg=white,bg=hitblue}
\setbeamercolor{secondary}{fg=hitblue,bg=hitlight}
\setbeamercolor{structure}{fg=hitblue}
\setbeamercolor{alerted text}{fg=hitgold}

% ─────────────────────────────────────────────────────────────────────────────
% PACKAGES
% ─────────────────────────────────────────────────────────────────────────────
\usepackage{amsmath}
\usepackage{amssymb}
\usepackage{mathtools}
\usepackage{booktabs}
\usepackage{graphicx}
\usepackage{hyperref}
\usepackage{tikz}
\usepackage{pgfplots}
\usepackage{tcolorbox}
\usepackage{listings}
\usepackage{xcolor}
\usepackage{algorithm}
\usepackage{algpseudocode}

% ─────────────────────────────────────────────────────────────────────────────
% TikZ LIBRARIES
% ─────────────────────────────────────────────────────────────────────────────
\usetikzlibrary{shapes,arrows,positioning,calc,chains,scopes}
\pgfplotsset{compat=1.17}

% ─────────────────────────────────────────────────────────────────────────────
% CODE LISTINGS
% ─────────────────────────────────────────────────────────────────────────────
\lstset{
  basicstyle=\ttfamily\small,
  breaklines=true,
  commentstyle=\color{gray},
  keywordstyle=\color{hitblue},
  stringstyle=\color{hitgold},
  showstringspaces=false,
  backgroundcolor=\color{hitlight},
  frame=single,
  rulecolor=\color{hitblue}
}

% ─────────────────────────────────────────────────────────────────────────────
% TCOLORBOX STYLES
% ─────────────────────────────────────────────────────────────────────────────
\tcbset{
  hitbox/.style={
    colback=hitlight,
    colframe=hitblue,
    fonttitle=\bfseries,
    boxrule=1pt
  },
  alertbox/.style={
    colback=yellow!10,
    colframe=hitgold,
    fonttitle=\bfseries,
    boxrule=1.5pt
  }
}

% ─────────────────────────────────────────────────────────────────────────────
% CUSTOM COMMANDS
% ─────────────────────────────────────────────────────────────────────────────

% Placeholder for figures
\newcommand{\placeholder}[3]{
  \begin{center}
    \fbox{\begin{minipage}[c][#2][c]{#1}
      \centering\small\color{gray}[Figure: #3]
    \end{minipage}}
  \end{center}
}

% Section divider frame
\newcommand{\sectionframe}[2]{
  \begin{frame}[plain]
    \begin{tikzpicture}[remember picture,overlay]
      \fill[hitblue] (current page.south west) rectangle (current page.north east);
      \node[anchor=center, white, font=\Large\bfseries] at (current page.center) {
        \begin{tabular}{c}
          #1 \\[0.3cm]
          {\small\color{hitgold}#2}
        \end{tabular}
      };
    \end{tikzpicture}
  \end{frame}
}

% ─────────────────────────────────────────────────────────────────────────────
% LOAD CUSTOM MACROS
% ─────────────────────────────────────────────────────────────────────────────
% ═══════════════════════════════════════════════════════════════════════════════
% SHARED MACROS — Mathematical Notation for Foundation Models Course
% ═══════════════════════════════════════════════════════════════════════════════

% ─────────────────────────────────────────────────────────────────────────────
% ACTIVATION & LAYER FUNCTIONS
% ─────────────────────────────────────────────────────────────────────────────
\newcommand{\softmax}{\mathrm{softmax}}
\newcommand{\sigmoid}{\mathrm{sigmoid}}
\newcommand{\relu}{\mathrm{ReLU}}
\newcommand{\gelu}{\mathrm{GELU}}
\newcommand{\LayerNorm}{\mathrm{LayerNorm}}
\newcommand{\FFN}{\mathrm{FFN}}
\newcommand{\MHA}{\mathrm{MHA}}
\newcommand{\Attn}{\mathrm{Attention}}

% ─────────────────────────────────────────────────────────────────────────────
% ATTENTION COMPONENTS
% ─────────────────────────────────────────────────────────────────────────────
\newcommand{\query}{Q}
\newcommand{\key}{K}
\newcommand{\val}{V}
\newcommand{\dmodel}{d_{\text{model}}}
\newcommand{\dk}{d_k}
\newcommand{\nheads}{h}

% Scaled dot-product attention formula
\newcommand{\SDPA}{
  \Attn(\query, \key, \val) = \softmax\left(\frac{\query \key^T}{\sqrt{\dk}}\right) \val
}

\newcommand{\CrossAttn}{\text{CrossAttention}}

% ─────────────────────────────────────────────────────────────────────────────
% PROBABILITY & EXPECTATION
% ─────────────────────────────────────────────────────────────────────────────
\newcommand{\E}{\mathbb{E}}
\newcommand{\Prob}{\mathbb{P}}
\newcommand{\KL}{\mathrm{KL}}
\newcommand{\ELBO}{\mathrm{ELBO}}
\newcommand{\given}{\mid}

% ─────────────────────────────────────────────────────────────────────────────
% NORMS & OPERATIONS
% ─────────────────────────────────────────────────────────────────────────────
\newcommand{\norm}[1]{\left\lVert #1 \right\rVert}
\newcommand{\abs}[1]{\left| #1 \right|}
\newcommand{\inner}[2]{\langle #1, #2 \rangle}
\newcommand{\T}{^{\top}}

% ─────────────────────────────────────────────────────────────────────────────
% VECTORS & MATRICES
% ─────────────────────────────────────────────────────────────────────────────
\newcommand{\bvec}[1]{\mathbf{#1}}
\newcommand{\mat}[1]{\mathbf{#1}}
\newcommand{\iid}{\stackrel{\text{i.i.d.}}{\sim}}
\newcommand{\defeq}{\coloneq}

% ─────────────────────────────────────────────────────────────────────────────
% LANGUAGE MODELING
% ─────────────────────────────────────────────────────────────────────────────
\newcommand{\vocab}{\mathcal{V}}
\newcommand{\context}{\text{context}}
\newcommand{\token}{t}
\newcommand{\seq}{x}
\newcommand{\ptheta}{p_\theta}
\newcommand{\pref}{p_{\text{ref}}}

% ─────────────────────────────────────────────────────────────────────────────
% RLHF & DPO
% ─────────────────────────────────────────────────────────────────────────────
\newcommand{\piref}{\pi_{\text{ref}}}
\newcommand{\pitheta}{\pi_\theta}
\newcommand{\yw}{y_w}
\newcommand{\yl}{y_l}
\newcommand{\DPOloss}{\mathcal{L}_{\mathrm{DPO}}}

% ─────────────────────────────────────────────────────────────────────────────
% RETRIEVAL & RAG
% ─────────────────────────────────────────────────────────────────────────────
\newcommand{\docs}{\mathcal{D}}
\newcommand{\qvec}{q}
\newcommand{\dvec}{d}

% ─────────────────────────────────────────────────────────────────────────────
% AGENTS & REASONING
% ─────────────────────────────────────────────────────────────────────────────
\newcommand{\obs}{o}
\newcommand{\act}{a}
\newcommand{\thought}{t}
\newcommand{\tools}{\mathcal{T}}
\newcommand{\history}{h}
\newcommand{\concat}{\oplus}

% ─────────────────────────────────────────────────────────────────────────────
% SETS & LOGIC
% ─────────────────────────────────────────────────────────────────────────────
\newcommand{\R}{\mathbb{R}}
\newcommand{\N}{\mathbb{N}}
\newcommand{\Z}{\mathbb{Z}}

\endinput


% ─────────────────────────────────────────────────────────────────────────────
% TITLE & AUTHOR (can be overridden in individual lectures)
% ─────────────────────────────────────────────────────────────────────────────
\author[D. Lederman]{Dror Lederman}
\institute{Holon Institute of Technology (HIT) \\ Data Science Department}
\date{\today}

% ─────────────────────────────────────────────────────────────────────────────
% FOOTER & HEADER
% ─────────────────────────────────────────────────────────────────────────────
\setbeamertemplate{footline}{
  \leavevmode%
  \hbox{%
    \begin{beamercolorbox}[wd=0.33\paperwidth,ht=2.25ex,dp=1ex,center]{author in head/foot}%
      \usebeamerfont{author in head/foot}\insertshortauthor
    \end{beamercolorbox}%
    \begin{beamercolorbox}[wd=0.34\paperwidth,ht=2.25ex,dp=1ex,center]{title in head/foot}%
      \usebeamerfont{title in head/foot}\insertshorttitle
    \end{beamercolorbox}%
    \begin{beamercolorbox}[wd=0.33\paperwidth,ht=2.25ex,dp=1ex,right]{frame number}%
      \usebeamerfont{frame number}\insertframenumber{} / \inserttotalframenumber\hspace*{2ex}
    \end{beamercolorbox}%
  }%
  \vskip0pt%
}

\endinput


\title{Week 11: Evaluation of Foundation Models and Agents}
\subtitle{Foundation Models and Agentic AI}
\author{Lecturer Name}
\institute{Holon Institute of Technology (HIT)}
\date{Semester, 2025}

\begin{document}

\begin{frame}
  \titlepage
\end{frame}

\begin{frame}{Learning Objectives}
  After this lecture you will be able to:
  \begin{itemize}
    \item Critically interpret standard \textbf{LLM benchmarks} (MMLU, HumanEval, BIG-Bench).
    \item Explain \textbf{hallucination evaluation} frameworks such as FActScore.
    \item Evaluate a \textbf{RAG system} using RAGAS and related metrics.
    \item Assess \textbf{agent reliability} via task success rate and error analysis.
    \item Design \textbf{human evaluation} protocols that are reproducible and unbiased.
    \item Discuss \textbf{cost, latency, and robustness} as first-class evaluation dimensions.
  \end{itemize}
\end{frame}

\section{Benchmarks for LLMs}
\sectionframe{LLM Benchmarks}{How do we measure what models know and can do?}

\begin{frame}{Standard Academic Benchmarks}
  \begin{center}
  \begin{tabular}{llll}
    \toprule
    \textbf{Benchmark} & \textbf{Domain} & \textbf{Format} & \textbf{Notes} \\
    \midrule
    MMLU~\cite{hendrycks2020mmlu} & 57 subjects & 4-way MCQ & Knowledge breadth \\
    HumanEval~\cite{chen2021humaneval} & Python coding & Pass@k & 164 problems \\
    BIG-Bench~\cite{srivastava2022bigbench} & 200+ tasks & Mixed & Emergent abilities \\
    GSM8K & Elementary math & Free-form & CoT eval standard \\
    MATH & Competition math & Free-form & Hard, 12K problems \\
    HellaSwag & Common sense & MCQ & Sentence completion \\
    TruthfulQA & Factuality & MCQ & Adversarial truths \\
    \bottomrule
  \end{tabular}
  \end{center}
\end{frame}

\begin{frame}{Benchmark Contamination and Gaming}
  \begin{tcolorbox}[alertbox, title=Critical Problem: Benchmark Contamination]
    Many training corpora contain benchmark test sets.
    A model may \textbf{memorize} test answers rather than demonstrating generalization.
  \end{tcolorbox}
  \vspace{0.4em}
  \textbf{Detection methods:}
  \begin{itemize}
    \item N-gram overlap between training data and benchmark.
    \item Canary strings inserted in private test sets.
    \item Dynamic benchmarks with new questions each evaluation round.
  \end{itemize}
  \vspace{0.3em}
  \textbf{Gaming:} Even without contamination, leaderboard pressure leads labs to
  over-optimize for public benchmarks. MMLU accuracy has limited correlation
  with real-world task performance.
  \note{Students should be skeptical of benchmark numbers. Always ask: was the model trained on this test set, even indirectly?}
\end{frame}

\begin{frame}{Pass@k for Code Evaluation}
  \textbf{HumanEval metric:} Pass@$k$ — the probability that at least one of $k$ samples passes.
  \vspace{0.4em}
  \begin{tcolorbox}[hitbox, title=Unbiased Estimator (Chen et al.\ 2021)]
    Given $n$ samples, $c$ of which pass:
    \[
      \text{Pass@}k = 1 - \frac{\binom{n-c}{k}}{\binom{n}{k}}
    \]
    Recommended: generate $n=200$, report Pass@1, Pass@10, Pass@100.
  \end{tcolorbox}
  \vspace{0.3em}
  \textbf{Limitation:} Unit tests are often weak — a function can produce wrong output on edge
  cases while passing the provided tests.
\end{frame}

\section{Hallucination Evaluation}
\sectionframe{Hallucination Evaluation}{Measuring factual faithfulness}

\begin{frame}{Types of Hallucination}
  \begin{columns}[T]
    \column{0.52\textwidth}
      \textbf{Intrinsic hallucination:}\\
      Output contradicts the provided context.
      \begin{itemize}
        \item ``The paper says X'' when the paper says Y.
      \end{itemize}
      \vspace{0.5em}
      \textbf{Extrinsic hallucination:}\\
      Output cannot be verified from any source.
      \begin{itemize}
        \item Fabricated citations.
        \item Invented statistics.
        \item Non-existent people or events.
      \end{itemize}
    \column{0.44\textwidth}
      \textbf{Measurement challenges:}
      \begin{itemize}
        \item Ground-truth often unavailable.
        \item Human annotation is expensive.
        \item LLM-as-judge has its own biases.
        \item Hallucination rates vary by topic and model.
      \end{itemize}
  \end{columns}
\end{frame}

\begin{frame}{FActScore: Fine-Grained Atomic Evaluation}
  \textbf{Min et al.\ (2023)}~\cite{min2023factscore}
  \vspace{0.4em}
  \begin{tcolorbox}[hitbox, title=FActScore]
    \begin{enumerate}
      \item Decompose long-form output into \textbf{atomic claims} (one fact per sentence).
      \item For each claim: retrieve relevant Wikipedia paragraphs.
      \item Verify: is the claim supported by the retrieved evidence?
    \end{enumerate}
    \[
      \text{FActScore} = \frac{\text{\# supported atomic claims}}{\text{\# total atomic claims}}
    \]
  \end{tcolorbox}
  \vspace{0.3em}
  \textbf{Findings:} GPT-4 hallucination rate $\approx$ 20\% on person biographies.
  Smaller models: 40–60\% hallucination rate.
\end{frame}

\section{RAG Evaluation}
\sectionframe{RAG Evaluation}{Measuring retrieval and generation quality jointly}

\begin{frame}{RAGAS Metrics}
  \textbf{RAGAS framework}~\cite{es2023ragas}:
  \vspace{0.3em}
  \begin{columns}[T]
    \column{0.50\textwidth}
      \begin{tcolorbox}[hitbox, title=Retrieval Metrics]
        \textbf{Context Precision:}\\
        $P = \frac{\text{relevant retrieved chunks}}{\text{total retrieved chunks}}$\\[4pt]
        \textbf{Context Recall:}\\
        $R = \frac{\text{relevant facts in context}}{\text{total relevant facts}}$
      \end{tcolorbox}
    \column{0.46\textwidth}
      \begin{tcolorbox}[hitbox, title=Generation Metrics]
        \textbf{Faithfulness:}\\
        Fraction of answer claims supported by context.\\[4pt]
        \textbf{Answer Relevance:}\\
        Similarity between the answer and reverse-generated questions.
      \end{tcolorbox}
  \end{columns}
  \vspace{0.3em}
  All metrics use LLM-as-judge — reference-free, scalable, but subject to judge bias.
\end{frame}

\begin{frame}{RAG Evaluation Checklist}
  \begin{tcolorbox}[hitbox]
    \begin{enumerate}
      \item \textbf{Retrieval quality:} MRR@k, NDCG@k on a labeled retrieval test set.
      \item \textbf{Context faithfulness:} \% of claims supported by retrieved context.
      \item \textbf{Answer correctness:} exact match / F1 against reference answers.
      \item \textbf{End-to-end latency:} retrieval + reranking + generation time.
      \item \textbf{Robustness:} performance on out-of-distribution or adversarial queries.
      \item \textbf{Cost:} embedding API calls + reranker + LLM tokens per query.
    \end{enumerate}
  \end{tcolorbox}
\end{frame}

\section{Agent Evaluation}
\sectionframe{Agent Evaluation}{Measuring task success and reliability}

\begin{frame}{Agent Benchmarks}
  \begin{center}
  \begin{tabular}{lll}
    \toprule
    \textbf{Benchmark} & \textbf{Domain} & \textbf{Metric} \\
    \midrule
    AlfWorld & Household tasks & Task success rate \\
    WebShop & Online shopping & Match score \\
    SWE-Bench & GitHub issues & \% issues resolved \\
    AgentBench & 8 diverse tasks & Normalized score \\
    GAIA & Real-world QA & Pass rate \\
    AppAgent & Mobile UI & Task completion \\
    \bottomrule
  \end{tabular}
  \end{center}
  \vspace{0.4em}
  \textbf{Key metric:} Task Success Rate (TSR) — \% of tasks fully completed without errors.
  \textbf{Secondary:} steps to completion, cost (tokens), failure mode distribution.
\end{frame}

\begin{frame}{Agent Failure Mode Analysis}
  \begin{tcolorbox}[alertbox, title=Common Agent Failure Modes]
    \begin{itemize}
      \item \textbf{Hallucinated tool call}: calls a tool with wrong arguments or non-existent tool.
      \item \textbf{Plan drift}: abandons the original task and pursues a tangent.
      \item \textbf{Infinite loop}: retries the same failing action repeatedly.
      \item \textbf{Partial success}: completes 4 of 5 sub-tasks but misses the key one.
      \item \textbf{Cascading errors}: one bad step propagates through the entire pipeline.
    \end{itemize}
  \end{tcolorbox}
  \vspace{0.3em}
  \textbf{Systematic evaluation:} log all agent traces; classify failures by type; prioritize the most frequent.
\end{frame}

\section{Human Evaluation}
\sectionframe{Human Evaluation}{The gold standard, and its pitfalls}

\begin{frame}{Human Evaluation Best Practices}
  \begin{columns}[T]
    \column{0.52\textwidth}
      \textbf{Setup considerations:}
      \begin{itemize}
        \item Blind/double-blind to prevent labeler bias.
        \item Clear rubric with 4–5 point Likert scale.
        \item Multiple raters per item; measure inter-rater agreement (Cohen's $\kappa$).
        \item Representative sample (not cherry-picked prompts).
        \item Include failure cases in the test set.
      \end{itemize}
    \column{0.44\textwidth}
      \textbf{Common dimensions:}
      \begin{itemize}
        \item \textbf{Helpfulness}: does it answer the question?
        \item \textbf{Accuracy}: factually correct?
        \item \textbf{Harmlessness}: safe, non-toxic?
        \item \textbf{Fluency}: natural, well-written?
        \item \textbf{Coherence}: internally consistent?
      \end{itemize}
  \end{columns}
\end{frame}

\section{Lab / Demo}
\begin{frame}{Lab Idea: Build an Evaluation Pipeline}
  \begin{tcolorbox}[hitbox, title=Lab 11 — LLM Evaluation Suite]
    \textbf{Goal:} Compare two LLMs on a custom evaluation suite.\\[4pt]
    \textbf{Tools:} \texttt{lm-evaluation-harness}, RAGAS, Python
    \begin{enumerate}
      \item Select 3 benchmark tasks (e.g., MMLU 100Q, GSM8K 50Q, HumanEval 20Q).
      \item Evaluate two models (e.g., Llama-3.2-1B vs.\ Llama-3.2-3B).
      \item Measure accuracy, latency, and estimated cost per query.
      \item Add a hallucination probe: 20 biography questions; score with FActScore (simplified).
      \item Write a 2-page evaluation report with findings and caveats.
    \end{enumerate}
    \textbf{Deliverable:} Evaluation report + code.
  \end{tcolorbox}
\end{frame}

\section{Discussion \& Summary}
\begin{frame}{Discussion Questions}
  \begin{enumerate}
    \item MMLU accuracy has become a standard leaderboard metric. A model trained with benchmark-contaminated data achieves 90\% MMLU but fails at real-world tasks. What does this say about the value of MMLU?
    \item FActScore decomposes responses into atomic claims and verifies them. What happens when the ``retrieval'' step used for verification is itself unreliable?
    \item Human evaluation is the gold standard, but it is expensive and slow. What are the conditions under which LLM-as-judge is a safe substitute for human evaluation?
    \item SWE-Bench requires resolving real GitHub issues. Is this the right benchmark for coding agents? What benchmarks would you design to measure agent reliability in production?
  \end{enumerate}
\end{frame}

\begin{frame}{Summary}
  \begin{itemize}
    \item Standard benchmarks (MMLU, HumanEval, BIG-Bench) measure capability but are vulnerable to \textbf{contamination} and \textbf{gaming}.
    \item \textbf{FActScore} enables fine-grained, scalable hallucination evaluation via atomic claim verification.
    \item \textbf{RAGAS} provides reference-free RAG evaluation on context precision, recall, faithfulness, and answer relevance.
    \item Agent evaluation requires \textbf{task success rate} + failure mode analysis; benchmarks like SWE-Bench, GAIA, and AgentBench are standard.
    \item \textbf{Human evaluation} is the gold standard; requires blinding, rubrics, and IAA measurement.
    \item \textbf{Cost and latency} are first-class evaluation dimensions alongside accuracy.
  \end{itemize}
  \vspace{0.3em}
  \textbf{Next week:} Safety, security, and alignment — what can go wrong and how to prevent it.
\end{frame}

\begin{frame}[allowframebreaks]{Suggested Reading}
  \nocite{hendrycks2020mmlu, chen2021humaneval, min2023factscore,
          srivastava2022bigbench, es2023ragas}
  \bibliography{../../references}
\end{frame}

\end{document}
